\documentclass[11pt]{article}
\usepackage[paperheight=11in, paperwidth=8.5in, margin=0.75in]{geometry}

\title{CALCULUS NOTATIONS}
\author{PHILIP THEURI}
\date{8$^{TH}$ June 2026}
\begin{document}
	\maketitle
	\section{RANGE AND DOMAIN}
The function $f(x)=(x-3)^2+\frac{1}{2}$ has a $\mathrm{D}_f:(-\infty,\infty)$ and range $\mathrm{R}_f:\left[\frac{1}{2},\infty\right)$.
	
The backslash left and right is to make sure that the brackets cover the whole fraction as require.
\section{LIMITS}
$\lim\limits_{x\to a}$\\[6pt]
$\lim_{x \to a}$\\

$\displaystyle{\lim\limits_{x\to a^-}}$
	
	
	$\displaystyle{\lim\limits_{x \to a }\, \, \, \, \frac{f(x)-f(a)}{x-a}=f(a)}$
\section{integrals}	
	
	$\int\sin x$\\
	$\int\cos x \,\mathrm{dx}$\\
	$\int \sin x \, dx=-\int \cos x + C$
	
	since the computer tries to fit the integral symbol into one line using the backslash display style is preferred so as the integral symbol can fit perfectly and cover the whole equation.\\
	$\displaystyle{\int \sin x \, dx =-\int \cos x + C}$\\
for definite integral ie: moving from a to b , we can use;\\	
	$\displaystyle{\int_a^b}$\\
	
	or use\\ $\displaystyle{\int\limits_a^b}$
	
	for definite integral you may encounter cases where you are required to move from say 2a to but that won't work the normal way hence use ;
	$\displaystyle{\int\limits_{2a}^{b_1}}$
	
	$W=\displaystyle{\int\limits_{p_1}^{p_2} \, \, \,-\left[V\right]}$
	
	$\displaystyle{\int\limits_{a}^{b}x^2 \, \mathrm{dx}=\left[\frac{x^3}{3}\right]_{a}^{b}=\,  \left(\frac{b^3}{3}\right)-\, \left(\frac{a^3}{3}\right)}$\\
	
	
	\section{\underline{Summation notation}.}
	
$\displaystyle{\sum\limits_{n=1}^{\infty}ar^n =a+ar+ar^2+\ldots+ ar^n}$\\[6pt]	

$\displaystyle{\int\limits_{a}^{b}f(x) \, \mathrm{dx}=}$


$\displaystyle{\int_a^b f(x) \, \mathrm{dx}=\lim\limits_{x \to \infty}\,\sum\limits_{k=1}^{n} f(x_k)\cdot \Delta x }$

the vector symbol is backslash vec\\
$\vec[v]=v_1\vec[i]+v_2\vec[j]=\langle v_1, v_2 \rangle$

\end{document}